\bigskip
\begin{center}\textls[120]{d) Grundlagen der Analysis. Komplexe Grössen.}\end{center}
\indent \textsc{Gauss’} freundschaftliche Beziehungen zu \textsc{Pfaff} begannen im Oktober 1798.
Mit dem Schluss des Sommersemesters 1798 hatte \textsc{Gauss} sein akademisches
Triennium beendet und war, wie er an \textsc{Bolyai} schreibt\footnote{\textit{Briefwechsel Gauss-Bolyai}, S. 5; die Briefe von \textsc{Gauss} an den 1798 noch in Göttingen zurückgebliebenen Freund, sind biographisch ausserordentlich interessant und wertvoll.}, am Dienstag, den
25. September nach Braunschweig zurückgekehrt. Damit hörte auch die ihm
vom Herzog \textsc{Carl Wilhelm Ferdinand} gewährte Studienunterstützung auf und
\textsc{Gauss} wollte sich durch Unterricht etwas zu verdienen suchen; über darauf
hinzielende Versuche berichtet er am 30. September und am 29. November
an \textsc{Bolyai}, sie führten nicht zum Ziel und er »lebt jetzt grossen Teils auf\par
\clearpage\SourceLeaf{130}\ArithmeticHead{45}{R}{III d. GRUNDLAGEN DER ANALYSIS. KOMPLEXE GRÖSSEN.}
\noindent Kredit«. Wie \textsc{Sartorius v. Waltershausen} berichtet, hat sich \textsc{Gauss} noch ein
Jahr vor seinem Tode (1854) wie folgt geäussert: »Ich habe in frühern Zeiten
wohl daran gedacht, Unterricht in der Mathematik geben zu müssen und ich
hatte mir zu diesem Ende ein Papier ausgearbeitet, das ich noch vor einigen
Jahren gesehen habe, das aber vielleicht nicht mehr existiert. Ich hatte darin
meine Gedanken über die Metaphysik der Mathematik niedergelegt .\ .\ .« Dieses
Papier hat sich im Nachlass gefunden und ist Werke XII, S. 57 ff. abgedruckt,
wir kommen nachher auf seinen Inhalt zurück. — Da es mit Unterricht in
Braunschweig nichts wurde, konzentrierte \textsc{Gauss} seine Tätigkeit auf die
Fertigstellung der \textit{Disquisitiones arithmeticae} und ging im Oktober 1798 nach
Helmstedt, um die dortige Bibliothek zu benutzen. »In Helmstedt bin
ich gewesen« — schreibt er an \textsc{Bolyai} — »und habe sowohl bei \textsc{Pfaff}
als bei dem Aufseher der Bibliothek eine sehr gute Aufnahme gefunden.
\textsc{Pfaff} hat meinen Erwartungen entsprochen. Er zeigt das untrügliche Kennzeichen
des Genies, eine Materie nicht eher zu verlassen, als bis er sie womöglich
ergrübelt hat.« Bald nach seiner Rückkehr in die Heimat gelang
es, vom Herzoge eine Fortsetzung der Unterstützung zu erlangen, er berichtet
am 30. Dezember an \textsc{Bolyai}: »In meiner Lage sind seit meinem letzten Briefe
einige günstige Veränderungen vorgegangen; .\ . der Herzog .\ .\ . hat erklärt,
dass ich die Summe, die ich in Göttingen genossen hatte, auch künftig behalten
solle (welche sich auf 158 Taler beläuft jährlich .\ .\ .\ .). Er wünscht
ferner, dass ich Dr. der Philosophie werde, ich werde es aber solange aufschieben,
bis mein Werk fertig ist, wo ich es hoffentlich ohne Kosten und
ohne die gewöhnliche Harlequinerie werde werden können.« Da sich aber der
Druck der \textit{Disqu.} zu lange hinzog, so entschloss er sich »eine andere Pièce
von einigen Bogen abdrucken zu lassen, die er nach Helmstedt schicken werde,
um .\ . die Dr.-Würde zu erwerben« (an \textsc{Bolyai}, am 22. April 1799). Es war
der Beweis der Wurzelexistenz für eine algebraische Gleichung, den er nach
Angabe des \textit{Tagebuchs} (Nr. 80) im Oktober 1797 gefunden hatte. Im Mai
sandte er diese Arbeit an \textsc{Pfaff}, zwei Briefe von diesem vom 30. Mai und
8. Juli 1799 (Werke X, 1, S. 99—105) zeigen, mit welchem Interesse und
warmem Verständnis \textsc{Pfaff} dieses Erstlingswerk aufgenommen hat. Unter
dem 26. Juni richtet \textsc{Gauss} an den Dekan, den Professor der Philosophie\par
\clearpage\SourceLeaf{131}\ArithmeticHead{46}{L}{L. SCHLESINGER, ÜBER GAUSS’ FUNKTIONENTHEORETISCHE ARBEITEN.}
\noindent \textsc{G. E. Schulze}, sein lateinisches Promotionsgesuch\footnote{Dieses ist im Original und in deutscher Übersetzung abgedruckt in der Nr. 15. des Braunschweigischen Magazins vom 16. Juli 1899 in einem Aufsatze von Dr. \textsc{Paul Zimmermann}, der überhaupt alle Akten betreffend das Promotionsverfahren wiedergibt.} aus dem wir einige charakteristische
Stellen hierhersetzen:

.\ .\ .\ . »Obwohl ich schon von frühester Kindheit an von leidenschaftlicher
Liebe zu den Wissenschaften, insbesondere zu den mathematischen,
beseelt war und keine Gelegenheit vorübergehen liess, wo ich
mein heisses Bestreben befriedigen konnte — so weit wenigstens wie ich,
von fast allen Hülfsmitteln entblösst, dies zu tun vermochte: — so wäre
dennoch wegen meiner Familienverhältnisse nur geringe Hoffnung gewesen,
das Ziel meiner Sehnsucht so, wie ich wünschte, erreichen und mich ganz
den Wissenschaften hingeben zu können, wenn ich nicht durch ein glückliches
Geschick unserm gnädigsten Herzoge bekannt geworden wäre. Nicht
nur ermuntert durch seine Huld und seinen Zuspruch, den Versuch zu machen,
mich den Wissenschaften zu widmen, sondern auch durch seine Freigebigkeit
auf das gnädigste unterstützt, habe ich zunächst auf den trefflichen Anstalten
unserer Stadt meine allgemeine Bildung gewonnen; dann aber habe ich, nachdem
ich jene Schulen durchlaufen, die Akademie zu Göttingen aufgesucht.
Dass ich diese der des Vaterlandes vorzog, hatte hauptsächlich seinen Grund
in dem Rufe der öffentlichen Bibliothek und der Fülle der literarischen Hülfsmittel,
ohne die, wie sattsam bekannt ist, Niemand auch nur Mässiges in den
mathematischen Fächern zu leisten imstande ist. Dort brachte ich mein
Triennium zu, meine ganze Zeit der Mathematik widmend, und hernach, in
dem nun abgelaufenen Jahre, kehrte ich hierher nach Hause zurück.

Durch so zahlreiche und bedeutende Wohltaten verpflichtet unserem besten
Fürsten, dessen Gnade mir bis auf diesen Tag niemals gefehlt hat, würde ich
es für ein Unrecht halten, wollte ich seinen Willen nicht als einen Befehl
und meine heiligste Pflicht ansehen. Deshalb habe ich, als er vor einiger Zeit
mir zu verstehen gab, es schiene ihm angemessen, dass ich die Würde eines
Doktors der Philosophie mir erwürbe, mich sofort angeschickt diesem Wunsche
nachzukommen; und da vor Allem erfordert wird, dass ein Kandidat durch
irgend eine Leistung seine Würdigkeit nachweist, so habe ich über ein analytisches
Thema von keineswegs geringer Bedeutung einen Aufsatz abgefasst\par
\clearpage\SourceLeaf{132}\ArithmeticHead{47}{R}{III d. GRUNDLAGEN DER ANALYSIS. KOMPLEXE GRÖSSEN.}
\noindent und schon vor einem Monate an den hochberühmten Herrn \textsc{Pfaff} abgesandt,
der ihn Ihnen vorlegen wird. Jetzt bitte ich daher Sie, ehrwürdigste Herren,
dass Sie nach Prüfung jener Abhandlung mir gemäss dem Rechte Ihrer Fakultät
die höchsten Ehren in der Philosophie übertragen.«

Schon am 28. Juni erstattete \textsc{Pfaff} sein Gutachten, in dem es heisst: »Ich
kann von dieser Abhandlung nicht anders als sehr vorteilhaft urteilen, da sie
von des Verfassers vorzüglichen Fähigkeiten und gründlichen Einsichten einen
überzeugenden Beweis enthält, so dass nach deren demnächst zu erwartenden
Abdrucke der Kandidat unter diejenigen zu rechnen sein wird, deren Promotion
unserer Fakultät zur Ehre gereicht.« Am 16. Juli erfolgte die Promotion
in absentia, noch ehe der Druck der Dissertation völlig abgeschlossen war.

Die beiden Ereignisse, der Abschluss der Dissertation und die Beobachtung
vom 30. Mai, stempeln den Mai 1799 zum Endpunkt der ersten
Jugendperiode, und wir haben nunmehr noch die in diese Periode fallenden
anderweitigen Untersuchungen von \textsc{Gauss} zu besprechen, soweit sie dem hier
zu schildernden Gebiet der Funktionentheorie angehören. Es sind dies einmal
die Beiträge zur Grundlagenforschung oder, wie \textsc{Gauss} es nennt, »Metaphysik
der Mathematik«, das andere Mal eben die Inauguraldissertation selbst.

\textsc{Gauss’} Interesse an den Grundlagen der Geometrie und Analysis geht
bis in die früheste Jugendzeit zurück, namentlich hat er sich nach seinen
eigenen Angaben (vergl. \textsc{Stäckel}, Abh. 4 dieses Bandes, Nr. 3, S. 17, 18)
mit den auf das \textsc{Euklidsche} Parallelenaxiom bezüglichen Fragen schon im Jahre
1792 beschäftigt, in welchem Jahre er auch den \textsc{Baermannschen} \textit{Euklid}\footnote{\textsc{G. F. Baermann}, \textit{Elementorum Euclidis libri XV}, Leipzig 1769.} sein
eigen nennen konnte. Aus der Göttinger Studentenzeit wissen wir, dass er
mit \textsc{Wolfgang Bolyai} viel über Grundlagen diskutiert hat, soll er doch,
nach dem Bericht von \textsc{Sartorius v. Waltershausen} (a. a. O. Fussnote\textsuperscript{13)},
S. 17) gesagt haben, \textsc{Bolyai} sei »der einzige gewesen, der in seine metaphysischen
Ansichten über Mathematik einzugehen verstanden habe.« Wenn
\textsc{Bolyai} (vergl. oben S. 23) sagt, \textsc{Gauss} sei damals in den Elementen der Arithmetik
und Geometrie weniger stark gewesen als er selbst, und wenn er später
(am 27. Dezember 1808) schreibt »schön wäre es, dass Du oben an den Spitzen
der stolzen Türme arbeitest, ich grübele an ihren Gründen«, so wird man\par
\clearpage\SourceLeaf{133}\ArithmeticHead{48}{L}{L. SCHLESINGER, ÜBER GAUSS’ FUNKTIONENTHEORETISCHE ARBEITEN.}
\noindent diesen Äusserungen einen gewissen Wahrheitsgehalt in bezug auf die mathematischen
Neigungen der beiden Freunde nicht absprechen können. Während
\textsc{Bolyai} die Grundlagenforschung an sich, um ihrer philosophischen Bedeutung
willen interessierte, dürfte \textsc{Gauss} von Anfang an mehr an konkreten Problemen
Gefallen gefunden, und in der Grundlagenforschung nur die solide Fundierung
für seine und anderer mathematische Spekulationen gesehen und gesucht haben.
Die Tatsache, dass ihm diese Fundierung für die Zahlentheorie am wenigsten
Schwierigkeiten bereitete, wird ihn mit dazu bestimmt haben, seine Bemühungen
zunächst dieser Disziplin zuzuwenden, und auch sein oft zitierter Ausspruch,
die Zahlentheorie sei die Königin der Mathematik, so wie diese selbst die
Königin unter den Wissenschaften, dürfte zum Teil aus solchen Erwägungen
hervorgegangen sein. Wenn er in der oben erwähnten Aufzeichnung (Werke
XII, S. 57) bei den arithmetischen Beziehungen von Grössen unterscheidet,
»ob der Begriff des Unendlichen vorausgesetzt werden muss oder nicht« und
sagt »der erste Fall gehöre in die höhere, der letztere in die gemeine oder
niedere Mathematik«, so hat er dabei sicher die Zahlentheorie in die zweite
Kategorie gezählt und für die höhere Mathematik die Schwierigkeit, die mit
dem Begriff des Unendlichen einhergeht, darin gesehen, dass dieser Begriff stets
nur im Sinne des Rigor antiquus, der bei den Alten, namentlich bei \textsc{Euklid}
und \textsc{Archimedes}, gewohnten Strenge, d. h. also gemäss der Exhaustionsmethode,
zur Geltung gebracht werden darf.

In diesem Sinne wird es auch zu verstehen sein, dass er bei dem Begriff
der reellen Zahl oder Grösse (die beiden Ausdrücke kommen abwechselnd bei
ihm vor) keine weiteren Schwierigkeiten empfunden haben dürfte, wie z. B.
daraus hervorgeht, dass er im Art. 16 der Dissertation mit einer gewissen
schlichten Selbstverständlichkeit sagt: »wir nehmen, um alle Strecken durch
Zahlen ausdrücken zu können, eine willkürliche Strecke als Einheit an«\footnote{Vergl. die Anfangsworte von \textsc{Newtons} \textit{Arithmetica universalis}, Cantabrigiae 1707, »Unter einer
abstrakten Zahl verstehen wir das Verhältnis einer beliebigen Grösse zu einer andern von derselben Art, die
als Einheit genommen wird«.},
und dass er in einer sogleich näher zu besprechenden Handschrift (Werke
X, 1, S. 391) eine veränderliche reelle Grösse »durch alle Zwischengrössen
hindurch stetig abnehmen« oder zunehmen lässt. Er stand eben in bezug auf
den Begriff der reellen, insbesondere der irrationalen Zahl ganz auf dem Stand-\par
\clearpage\SourceLeaf{134}\ArithmeticHead{49}{R}{III d. GRUNDLAGEN DER ANALYSIS. KOMPLEXE GRÖSSEN.}
\HistoricalFoot{X$_2$. Abh. 2.}{7}
\noindent punkt seiner Zeitgenossen und hat darin volle Befriedigung gefunden. Dagegen
scheint er in der Ausgestaltung der Begriffe, die mit einer unendlichen
Menge von reellen Zahlen einhergehen, schon sehr frühzeitig seine eigenen
Wege eingeschlagen zu haben. — Während die Werke XII, S. 57 abgedruckte
Handschrift dem Aufbau der elementaren Arithmetik gewidmet ist, scheint
es, dass wir in der Werke X, 1, S. 390 abgedruckten eine Art Fortsetzung
der ersteren in das Gebiet der »höheren Mathematik« vor uns haben. Die
deutsche Redaktion lässt darauf schliessen, dass \textsc{Gauss} auch diese für etwaige
Unterrichtszwecke entworfen haben mag, ihr Aufbau ist durchaus elementar
und geht vorsichtig vom einfacheren zum schwierigeren. In dieser Handschrift
betrachtet er »Reihen« von reellen Grössen, d. h. in der uns geläufigen
Terminologie, Folgen, also abzählbare Mengen, die in einer bestimmten
Anordnung vorliegen, und bemerkt, dass die Betrachtung »des Inbegriffs einer
jeden Anzahl solcher Grössen«, also, wie wir sagen würden, beliebiger linearer
Punktmengen, »wenig Nutzen haben würde«. Übrigens sind die von \textsc{Gauss}
angestellten Betrachtungen ohne weiteres auch auf beliebige lineare Punktmengen
anwendbar.

Eine solche »Reihe« kann nach oben bezw. unten hin beschränkt sein
und besitzt dann eine kleinste obere bezw. grösste untere Grenze oder, wie wir
lieber sagen wollen, \textls[120]{Schranke}, deren Existenz \textsc{Gauss} durch einen Stetigkeitsschluss
beweist. Durch Weglassen von Gliedern der Reihe kann die obere
Schranke nicht vergrössert, die untere nicht verkleinert werden; unterdrückt
man also sukzessive immer mehr Glieder der Reihe, so bilden die oberen
Schranken eine nicht zunehmende, die unteren eine nicht abnehmende Reihe,
und zwar ist, wenn Beschränktheit nach beiden Seiten vorliegt, jedes Glied
der Reihe der obern Schranken nicht kleiner als irgendein Glied der Reihe
der untern. — Die erste Reihe besitzt daher eine untere Schranke \(L\), die
zweite eine obere Schranke \(M\), und es ist stets \(L\geqq M\). — \textsc{Gauss} nennt \(L\)
die letzte obere, \(M\) die letzte untere Grenze, es sind offenbar die von \textsc{Paul du
Bois Reymond} sogenannten Unbestimmtheitsgrenzen:
\[
L=\lim\sup,\quad M=\lim\inf.
\]
Wenn für eine beiderseits beschränkte Reihe \(L=M\) gilt, so hat die Reihe
eine absolute Grenze, d. h. sie strebt nach einem bestimmten endlichen Grenz-\par
\clearpage\SourceLeaf{135}\ArithmeticHead{50}{L}{L. SCHLESINGER, ÜBER GAUSS’ FUNKTIONENTHEORETISCHE ARBEITEN.}
\noindent wert, einem Limes. — Es ist bei dieser Darstellung, die zur Zeit ihrer Publikation
(1912 im Heft III der \textit{Materialien}), trotzdem sie schon über hundert
Jahre alt war, durchaus neuartig erschien, besonders bemerkenswert, dass
\(\lim\sup\) und \(\lim\inf\) nur auf Grund der Anordnung, also ohne Abstandsbegriff
eingeführt werden. Der Begriff des Häufungspunktes, mit dem z. B. \textsc{Bolzano}
operiert, kommt bei \textsc{Gauss} überhaupt nicht vor, und \textsc{Gauss} scheint hierauf
besonderen Wert gelegt zu haben. Sagt er doch in einem aus späterer Zeit
(1825) stammenden Fragment (Werke X, 1, S. 396), das sich allerdings vorwiegend
auf komplexe Grössen bezieht, »es sei ganz vorzüglich wichtig, die
Theorie des Gegensatzes« — d. h. der Unterscheidung zwischen den verschiedenen
Elementen — »zur Klarheit zu bringen ohne \textls[120]{Grösse}«, was soviel besagen
will, wie ohne \textls[120]{Abstandsbegriff}.

Den Begriff der obern Schranke benutzt \textsc{Gauss} auch schon im Art. 6
der Dissertation als etwas ihm durchaus vertrautes; auch ist er mit seiner
scharfen Erfassung des Limesbegriffs ohne weiteres in der Lage, die Fragen
der Konvergenz und Divergenz unendlicher Reihen (Summen) zu behandeln.
In der Tat finden wir im Art. 6 der Dissertation, in der Fussnote zum Absatz 3,
eine Erörterung über Konvergenz und Divergenz von Potenzreihen,
auf die \textsc{Gauss} selbst bei Gelegenheit (1850, Werke X, 1, S. 435) hinweist,
und in dem Nachlassstück über die Summe der divergenten Reihe \(1-1+
2-6+\ldots\) (Werke X, 1, S. 382—384) wird von einer Potenzreihe gesagt, sie
konvergiere für \(\omega<1\), während die Konvergenz für \(\omega=1\) zwar benutzt, aber
nicht ausdrücklich bewiesen wird. In eben diesem Nachlassstück zeigt \textsc{Gauss}
auch, dass er dem Begriff der Summe gewisser divergenter Reihen einen präzisen
Sinn beizulegen weiss, und dass er mit der durch divergente Potenzreihen,
die linearen Differentialgleichungen genügen, vermittelten asymptotischen
Annäherung vertraut war, wobei an den Einfluss der Arbeiten von \textsc{Laplace}
(1785) zu denken sein wird. — Dass er die asymptotische Darstellung der
\textsc{Besselschen} Funktion \(J_1(x)\) für grosse Werte des Arguments gekannt hat, zeigt
die Tagebuchnotiz 82. vom 16. Oktober 1797, die (vergl. oben S. 26) einen
asymptotischen Ausdruck für die Nullstellen der gedachten Funktion gibt.
Näheres hierüber, sowie über die Untersuchung des asymptotischen Verhaltens
entfernter Glieder der Entwicklung der Mittelpunktsgleichung findet man in\par
\clearpage\SourceLeaf{136}\ArithmeticHead{51}{R}{III d. GRUNDLAGEN DER ANALYSIS. KOMPLEXE GRÖSSEN.}
\noindent den Bemerkungen Werke X, 1, S. 385—389 und S. 438—447 (vergl. auch unten
Abschn. V e)).
